%!TEX root = ./../../thesis-trier.tex

\index{$m$-cycle|see{cycle}} 
\IN{cycle} 
\index{cyclic permutation|see {cycle}} 
\IN{cyclically permuting a set}
\IN{transposition}
\IN{cycle notation} \label{def:cycle-notation}
\IN{disjoint cycles}
\IN{cycle induced by transformation}
\IN{cycle contained in transformation!see {cycle induced by transformation}}
\IN{permuted subset}
\IN{point contained in cycle}
Let $X$ be some set and $a_0, a_1, \ldots, a_{m-1} \in X$ be $m \ge 1$ distinct elements.
Then the permutation $(a_0 ~ a_1 ~ \cdots ~ a_{m-1})_X \colon X \to X$ 
given by 
\[
(a_0 ~ a_1 ~ \cdots ~ a_{m-1})_X(x) = 
\begin{cases} 
	a_{(j+1) \bmod m} & \text{if } x = a_j \text{ for some } j \in \{0,1,\ldots,m-1\}, \\
	x & \text{otherwise,}
\end{cases} 
\]
is called an \emph{$m$-cycle} or \emph{cycle} of length $m$.\footnote{Note that for $X \ne \{a_0,a_1, \ldots, a_m\}$
	the cycle has fixpoints. Some authors exclude this possibility and 
	stipulate $X = \{a_0,a_1, \ldots, a_m\}$.}

Equivalently, $t \colon X \to X$ is an $m$-cycle if and only if
$m \le 1$ and $t = \operatorname{id}_X$ 
or $m = |\move{t}|$ and for $x,x' \in \move{t}$
there exists $i \ge 0$
such that $t^i(x) = x'$.
When the context is clear, we often omit the index~$X$ and simply write $(a_0 ~ a_1 ~ \cdots ~ a_{m-1})$.
We do not use commas in the cycle notation between points, 
to distinguish it from an ordinary $m$-tuple.
When writing a $1 \times n$ matrix as in \cref{prelim:sync_max:matrixscheme} 
we obtain the same notation as for a cycle,
but in each case the context should make it clear whether we refer to a cycle or a $1 \times n$ matrix.
For each $a \in X$, we have $(a) = () = \operatorname{id}_X$.
The following equations hold:
\begingroup
\allowdisplaybreaks 
\begin{align}
	(a_0 ~ a_1  ~\cdots~  a_{m-2} ~ a_{m-1}) & = (a_1 ~ a_2 ~ \cdots ~ a_{m-1} ~ a_0) \\
	(a_0 ~ a_1  ~\cdots~  a_{m-1}) & = (a_1 ~\cdots~ a_{m-1})(a_0 ~ a_{m-1}) = (a_0 ~ a_2 ~ \cdots ~ a_{m-1}) (a_0 ~ a_1) \label{prelim:eqn:cycletrans1} \\
	(a_0 ~ a_1  ~\cdots~  a_{m-1}) & = (a_{m-2} ~ a_{m-1})(a_0 ~ a_1  ~\cdots~  a_{m-3} ~ a_{m-1}) \label{prelim:eqn:cycletrans2} \\		
	(a_0 ~ a_1  ~\cdots~  a_{m-1}) & = (a_0 ~ a_1)( a_1 ~\cdots~ a_{m-1}) \\		
	(a_0 ~ a_1  ~\cdots~ a_{m-2} ~ a_{m-1}) & = \concat{\concat{(a_0 ~\cdots~ a_{m-3})}{(a_{m-2} ~ a_{m-1})}}{(a_{m-2} ~ a_0)} 
\end{align}
\endgroup%
% einige oben nur spezialfälle davon
And for $0 \le i < j < m$ we have:
\begin{align*}
	(a_0 ~ a_1 ~\cdots~ a_{m-1})(a_i ~ a_j)
	& = (a_0 ~ a_1 ~\cdots~ a_{i} ~ a_{j+1} ~ a_{j+1} \cdots a_{m-1})(a_{i+1} ~\cdots ~ a_j) \\
	(a_i ~ a_j)(a_0 ~ a_1 ~\cdots~ a_{m-1})
	& = (a_0 ~ a_1 ~\cdots~ a_{i-1} ~ a_j ~ a_{j+1} ~\cdots~ a_m)(a_i ~ \cdots ~ a_{j-1})
\end{align*}

A \emph{transposition} is a $2$-cycle. Two cycles~$t,h \colon X \to X$ are \emph{disjoint} if $\move{t} \cap \move{h} = \emptyset$,
in this case we have $th = ht$.  
Let $t \colon X \to X$ be some transformation.
For finite $S \subseteq X$, we say that $t$
\emph{cyclically permutes $S$} if $t$ permutes $S$
and $t_{|S}$ is an $|S|$-cycle. 
Then $t^{|S|}(x) = x$ for all $x \in S$. A permutation on a set of size at least two cyclically permutes its domain
if and only if it is a cycle with no fixpoints.
If $t_{|S} = (a_0 ~ a_1 ~ \ldots ~ a_{m-1})$
we also say that $t$ 
\emph{induces the cycle}, or \emph{contains the cycle}, $(a_0 ~ a_1 ~ \cdots ~ a_{m-1})$.
Equivalently, $t(a_j) = a_{(j+1) \bmod m}$ for all $j \in \{0,1,\ldots,m-1\}$, i.e.,
$(x ~ t^2(x) ~ \cdots ~ t^{m-1}(x)) = (a_0 ~ a_1 ~ \ldots ~ a_{m-1})$.
%Then, for each $x \in \{a_0, a_1, \ldots, a_{m-1}\}$,
%we have $t^m(x) = x$. Conversely, 
For $x \in X$ and $m \ge 1$, the points 
$x, t(x), \ldots, t^{m-1}(x)$ form a cycle induced by $t$ if and only if~$m$ is minimal
such that $t^m(x) = x$.
Two distinct cycles induced by $t$ are disjoint.

\begin{proposition} \label{prelim:fact:product_disjoint_cycles} 
	A permutation on a finite set equals the composition (in any order) of its induced cycles.	
\end{proposition}
\begin{proof}
	Let~$t \colon [n]\to [n]$ be a permutation
	with induced cycles $t_1, \ldots, t_m$.
	For $x \in [n]$, we have $t_i = (x ~ t(x) \cdots t^{k-1}(x))$
	for some minimal $k \ge 1$ with $t^k(x) = x$ and 
	unique $i \in \{1,\ldots,m\}$.
	Furthermore, for $j \ne i$, we also have $t_j(x) = x$.
	Thus, the result of any composition of
	$t_1, \ldots, t_m$ maps $x$ to $t_i(x) = t(x)$,
	and so this composition equals~$t$.
\end{proof}

\begin{corollary} \label{prelim:fact:transpositions_generate_Sn}
	A permutation on a finite set can be written as
	a composition of transpositions.
\end{corollary}
\begin{proof}
	This follows from \cref{prelim:eqn:cycletrans1,prelim:fact:product_disjoint_cycles}.
\end{proof}